\renewcommand{\Anno}{2011}

\renewcommand{\NumeroProblema}{1}

\renewcommand{\Testo}{
Consideriamo la somma
$$\frac{1}{2}+\frac{1}{3}+\ldots+\frac{1}{11}+\frac{1}{12},$$
in cui compaiono tutte le frazioni con numeratore 1 e denominatore compreso tra 2 e 12.  Dalla somma si vogliono eliminare alcuni degli addendi, in modo che la somma dei restanti sia esattamente 1.

Determinare tutti i modi in cui si può fare questa operazione.
}

\renewcommand{\Risposta}{
Il problema ammette 2 soluzioni:
$$\frac{1}{2}+\frac{1}{3}+\frac{1}{6}=1,
\hspace{3em}
\frac{1}{2}+\frac{1}{4}+\frac{1}{6}+\frac{1}{12}=1.$$
}

\renewcommand{\Aiutino}{ 
È possibile che in una soluzione compaiano la frazione 1/7 o la frazione 1/11?  È possibile che in una soluzione compaiano la frazione 1/5 o la frazione 1/10?  È possibile che in una soluzione compaiano la frazione 1/9 o la frazione 1/8?  È possibile che in una soluzione compaia la frazione 1/4 ma non la frazione 1/12 (o viceversa)?  A questo punto non restano molti casi \ldots
}

\renewcommand{\SoluzioneUfficiale}{\TitoloSoluzione 
Il problema ha le seguenti 2 soluzioni:
$$\frac{1}{2}+\frac{1}{3}+\frac{1}{6}=1,
\hspace{3em}
\frac{1}{2}+\frac{1}{4}+\frac{1}{6}+\frac{1}{12}=1.$$

Resta da dimostrare che non ci sono altre possibilità, cosa che facciamo in 5 passi.

\subparagraph{\textmd{Passo 1}}

Mostriamo che in una soluzione del problema non possono comparire le frazioni con denominatore 7 e 11.

Supponiamo infatti che in una soluzione compaia la frazione 1/7. Allora la somma delle restanti frazioni deve essere uguale a 6/7. Questo non è possibile perché tutte le frazioni rimaste hanno un denominatore che non contiene fattori 7, quindi la loro somma non può avere il denominatore divisibile per 7.

Stesso discorso vale per il fattore 11.

\subparagraph{\textmd{Passo 2}}

Mostriamo che in una soluzione del problema non possono comparire le frazioni con denominatore 5 e 10.

Se infatti in una soluzione compare una sola delle 2, allora si ragiona come nel passo precedente considerando il fattore 5. Se invece sono presenti tutte e 2, allora la loro somma è 3/10, quindi la somma delle restanti deve essere 7/10, e ancora una volta questo non è possibile perché non ci sono più denominatori divisibili per 5.

\subparagraph{\textmd{Passo 3}}

Mostriamo che in una soluzione del problema non possono comparire le frazioni con denominatore 8 e 9.

Supponiamo infatti che in una soluzione compaia la frazione 1/9. Allora la somma delle restanti frazioni deve essere uguale a 8/9.  Le restanti frazioni però hanno al massimo un fattore 3 al denominatore, quindi il loro denominatore comune conterrà al massimo un fattore 3.

Il discorso per l'8 è analogo, basta contare i fattori 2 al denominatore.

\subparagraph{\textmd{Passo 4}}

Mostriamo che in una soluzione c'è la frazione 1/4 se e solo c'è la frazione 1/12.

Infatti già sappiamo che 1/8 non c'è, quindi per lo stesso ragionamento di prima non può esserci una soluzione che coinvolge una sola frazione con denominatore divisibile per 4.

\subparagraph{\textmd{Passo 5}}

Nei passi precedenti ci siamo ridotti a considerare le frazioni con denominatori 2, 3, 4, 6, 12.  Inoltre sappiamo che c'è 1/4 se e solo se c'è 1/12.  Ci restano quindi 2 casi.
\begin{itemize}
	\item  Se non ci sono 1/4 e 1/12, allora la somma delle frazioni restanti è esattamente 1, quindi abbiamo la prima soluzione.

	\item  Se ci sono 1/4 e 1/12, che come somma fanno 1/3, allora l'unica possibilità è di rinunciare alla frazione 1/3, ottenendo così la seconda soluzione.
\end{itemize}


}

\renewcommand{\SoluzioniAlternative}{}

\renewcommand{\Osservazioni}{}

\renewcommand{\Commenti}{
}